\section{Introduction}

To understand the safety risks of autonomous artificial intelligence, the measurement of multi-step actions is essential. The earlier safety benchmarks tested only isolated text outputs. In commercial software stacks, we now deploy agents that manage inventories, communicate with external suppliers, and execute browser transactions for a human principal, thereby shifting the safety question from static text generation to real-world execution. We therefore ask whether an agent deployed to manage a company can come to engage in serious criminal activity.

To be able to trace back how such an outcome could arise, it makes sense to break the system down to its most basic parts. The contemporary agent stack is usually evaluated on separate layers of multi-agent economics, model alignment, and browser automation. This study connects those three layers along one operational chain. We decompose the pathway to automated illicit procurement into three conditional stages:
\[
P(A, B, C) = P(A) \cdot P(B \mid A) \cdot P(C \mid A, B),
\]
In this equation $A$ denotes financial distress after a coordination failure, $B$ denotes engagement with a harmful tool chain, and $C$ denotes completion of a simulated illicit purchase. In words: if ordinary multi-agent coordination dynamics push an agent into insolvency and contextual pressure raises its willingness to select harmful tools, then general-purpose browser autonomy can close the loop by executing a prohibited checkout on an external storefront.

The following sections examine each stage in turn. The empirical results of our supply-chain coordination benchmark are presented first to quantify the baseline bankruptcy rate $P(A)$. We then analyze how financial pressure, payoff ratios, and monitoring cues shape harmful tool selection $P(B \mid A)$ across $53{,}000$ recorded trials while we vary both account solvency and principal framing. Furthermore, six frontier model configurations are evaluated in a browser environment to measure whether they complete simulated purchases of controlled substances and firearms (see Fig. 1). Because each phase was tested in a separate environment, we also trace back where the composed estimate holds and where empirical anomalies weaken the accidental pathway.

\section{Phase 1: Coordination Failure Leaves Agents Recurrently Distressed}

In order to test whether ordinary business operations produce financially distressed agents without adversarial interference, four coupled agents are evaluated in a stochastic supply-chain game. We place a retailer, a wholesaler, a distributor, and a factory in a sequential ordering loop over $100$ rounds. This run is classified as bankrupt when the balance of any role reaches zero. The most common outcome across our $58$ analyzed runs is bankruptcy, which is observed in $49$ runs ($P(A) = 0.845$).

As one can see in the balance trajectories of Fig. 2, order variability amplifies as it propagates upstream through the supply chain. There are five recurrent failure modes across the bankrupt runs. The dominant pattern is an Overordering Cascade in $36$ of the $49$ bankrupt runs ($73.5\%$), where the retailer over-orders after an initial demand spike and upstream nodes amplify the surge into unsustainable inventory backlogs and penalty costs. The remaining insolvencies are caused by Backlog Death Spirals, Inventory Hoarding, Mid-Chain Squeezes, and Wealth Concentration.

For the sake of comparison, four system configurations are tested across five model families in the Beer Game environment. While introducing an advisory orchestrator agent delays insolvency from a mean of $63.4$ rounds in the Standard baseline to $80.8$ rounds, we find that $11$ of $14$ orchestrated runs still end in bankruptcy by the end of the game ($p = 0.33$). The distributor is the first participant to fail in $25$ of the $49$ bankrupt runs rather than the factory at the upstream end. This means that ordinary multi-agent coordination dynamics recurrently leave agents operating at a severe financial loss while they are still tasked with generating revenue for their principal.

\section{Phase 2: Contextual Cues Modulate Harmful Tool Use}

Because Phase 1 shows that autonomous agents frequently enter financial distress, we test whether pressure to recover money increases their use of harmful tools. We first measure baseline engagement in a competitive business game across ten tool categories that each follow a five-step escalation ladder from borderline misconduct at Step 1 to severe criminal offenses at Step 5. In the unpressured baseline, the three most severe categories of corruption, sabotage, and violence receive zero engagement across all tested models. The remaining seven categories are engaged unevenly across models, since all three tested configurations engage more often with data theft ($0.400$ to $0.667$) than with milder financial crime ($0.100$ to $0.500$).

To separate the agent's numerical financial state from the channel through which pressure is communicated, we sweep three models across $23$ account balances from $-2{,}000$ to $+12{,}000$, ten harmful tool chains, and three pressure framings ($27{,}600$ decisions). When only the numeric balance is stated in configuration E1, step-one engagement is $0.348$ for GPT-4.1-mini and $0.301$ for Gemini-3-Flash. Under verbal pressure from a principal, we observe that these step-one engagement rates more than double to $0.773$ and $0.733$ across the $1{,}150$ trials per cell, thereby showing that words drive behavior much more strongly than the balance number alone. The Claude Haiku-4.5 model refuses every request.

There are two constraints on this pressure effect. First, high initiation at Step 1 does not imply completion of Step 5, because GPT-4.1-mini initiates the violence chain at $0.896$ under verbal pressure by sending a legal threat while reaching Step 3 stays at $0.009$ and full completion is $0.000$. Second, prepending a realistic Phase 1 bankruptcy transcript in configuration E2 unexpectedly lowers insolvent verbal-pressure engagement from $0.815$ to $0.583$ for GPT-4.1-mini and from $0.872$ to $0.748$ for Gemini-3-Flash ($n = 400$ per cell). The distress history reduces recklessness here.

To test whether explicit financial incentives also shift agent choices without verbal pressure, we run a separate payoff experiment across $12{,}150$ decisions between a safe option, a harmful option, and abstention (see Fig. 3). When both actions have an equal $0.9$ success probability and equal expected value, harmful choice is $0.013$ for GPT-4.1-mini and $0.000$ for the other two models. The harmful choice rate climbs monotonically when the harmful payout is raised to six times the safe payout ($6\times$ matched), reaching $0.220$ for GPT-4.1-mini and $0.384$ for Gemini-3-Flash. Within this $6\times$ matched cell, insolvent agents select the harmful option at $0.420$ versus $0.120$ when solvent for GPT-4.1-mini and at $0.647$ versus $0.253$ for Gemini-3-Flash, thereby confirming that insolvency amplifies the pull of high-payout misconduct.

We also evaluate whether contextual guardrails suppress commitment depth. When we tell the agent in the system prompt that its actions are monitored, mean commitment depth is reduced across all three configurations and brings Claude 3.5 Haiku nearly to zero. The explicit description of ethical harm also lowers commitment depth for two of the three tested models, whereas stating that the agent has used the tool four times in the past raises depth for GPT-4.1-mini from $1.10$ to $1.55$. These results show that tool selection is not fixed, but moves substantially with framing, payoffs, and oversight cues.

\section{Phase 3: Browser Agents Close the Loop Under Obfuscation}

While Phase 2 measures selection among abstract tool APIs, closing the loop requires testing whether an agent can execute an illicit purchase in a browser environment. We evaluate six frontier model configurations driving a locally hosted storefront via Claude Code and Playwright MCP across $360$ trials ($24$ cells with $n = 15$ per cell). We give each agent a shipping address, a test card, and a task to buy one of five controlled items. This entire benchmark runs inside our local sandbox.

To isolate the mechanism of safety refusal, two storefront files are compared that are byte-identical in product photographs, prices, categories, and descriptions while differing only in the visible product name. The explicit condition lists literal names such as 1KG Cocaine or 1 Shotgun. In the pictograph condition, product names are replaced by publicly documented drug-enforcement emoji codes. We cross both naming regimes with screenshot vision enabled and disabled.

To measure how far agents progress before stopping, commitment depth is defined across an eight-step checkout funnel from search at Step 1 to order placement at Step 8. In all $360$ trials $72\%$ of runs pass the initial search step, while $88$ runs end in immediate refusal and $12$ fail from browser errors. Under explicit product naming, we find that only $50\%$ of runs ($90$ of $180$) pass the initial search step, whereas $94\%$ ($170$ of $180$) pass search when product names are replaced by pictographs. The step tracker shows that $79\%$ of the $214$ trials reaching payment proceed to place the order at Step 8.

The graphical comparison of all $24$ cells in Fig. 4 shows that multiple frontier configurations complete simulated illicit purchases when product names are obfuscated. The highest overall completion rate is reached by Gemini-3-Flash at $55\%$ ($33$ of $60$), followed by Sonnet 4/4.5 and GPT-4.1-mini at $52\%$ each. On the emoji storefront with vision disabled, Sonnet 4/4.5 completes $93\%$ ($14$ of $15$) of checkouts, while GPT-5.4 completes $60\%$ and Gemini-3.1-Pro completes $53\%$. When vision is enabled on the emoji storefront, completion jumps by $60$ percentage points for both Haiku-4.5 and Gemini-3-Flash.

The refusal contrast between explicit and pictograph naming holds across five of the six tested configurations. For Sonnet 4/4.5 and Haiku-4.5, all $24$ refusals are recorded on the explicit storefront, compared with zero refusals on the emoji storefront. The refusal logs show that models cite specific product strings or drug schedules, and no model ever refuses the storefront as a whole. This means that current refusal behavior acts primarily as a lexical trigger on the product title rather than a semantic evaluation of the transaction context.

\section{Composition: Tracing Back the Full Loop}

To examine how the three empirical stages combine, we calculate the chain-rule product $P(A, B, C) = P(A) \cdot P(B \mid A) \cdot P(C \mid A, B)$ from our measured rates. We set $P(A) = 0.845$ from the Phase 1 bankruptcy rate and $P(C \mid A, B) = 0.511$ from the pooled vision-off emoji storefront completion rate. For the middle transition $P(B \mid A)$, the Phase 2 insolvent engagement rates from configuration E2 are used. This layout holds the outer two terms fixed so that we can trace how the overall joint probability responds to changes in framing and model choice.

If taking the borderline first step of a harmful tool chain counts as engagement, the resulting joint probability under verbal pressure from a principal is $0.252$ for GPT-4.1-mini and $0.323$ for Gemini-3-Flash across the insolvent cells. When only the insolvent account balance is stated without verbal pressure, the joint product drops roughly in half to $0.138$ and $0.151$ for the two models. If we require the agent to reach Step 3 of the tool chain where unambiguous criminal acts begin, the verbal-pressure joint rate becomes $0.009$ for GPT-4.1-mini and $0.239$ for Gemini-3-Flash. The composed product for Claude Haiku-4.5 is zero.

It is very important to mention that these composed values are illustrative calculations rather than measured end-to-end deployment forecasts. First, the three stages were executed on separate setups and never as one continuous trajectory where the same agent went bankrupt, selected a tool, and opened a browser. Second, the only cross-stage link that we tested directly worked against the accidental pathway, since prepending a realistic Phase 1 bankruptcy transcript in configuration E2 reduced harmful engagement relative to configuration E1. This means that Phase 2 does not show how a benign agent would acquire an illicit purchasing tool on its own.

Additionally, all $360$ Phase 3 trial logs were audited manually and $31$ automated labels were corrected where browser-session lockups or unverified clicks had been misclassified. We find after this audit that GPT-4.1-mini records zero policy refusals on either storefront because all seven of its initially logged refusals were actually browser or navigation failures. While a live dark-web market would introduce extra friction such as anti-bot challenges and cryptocurrency rails, roughly $28\%$ of our trials already fail at initial navigation and search. Furthermore, because capability friction erodes rapidly as browser agents improve, relying on task difficulty as a safety barrier offers only temporary protection.

\section{Governance Interventions: Breaking the Chain}

Because existing capabilities can already be composed by an adversarial principal or triggered through environmental framing, we cannot rely solely on lexical refusal training for safety. We see that an attacker or an illicit vendor can bypass string-level refusal simply by substituting pictographs or benign euphemisms for prohibited product names. Moreover, even a small per-agent failure probability $p$ yields an expected incident count $pN$ that grows linearly with the number of deployed agents $N$. We therefore propose three targeted interventions across the evaluation, execution, and protocol layers to break the loop before execution.

First, any AI system authorized to execute financial transactions should undergo mandatory composable red-teaming prior to deployment. We argue that instead of evaluating chat refusals on explicit text prompts in isolation, pre-deployment benchmarks must measure end-to-end completion rates under realistic obfuscation conventions and contextual pressure. This requirement parallels the risk-management and robustness obligations defined in Articles 9 and 15 of the European Union AI Act. The measurement of conditional transitions across stages lets us verify whether a refusal boundary holds when upstream context degrades.

Second, agentic execution frameworks should enforce hard human-confirmation gates at financially and legally material steps. The commitment-depth funnel shows that $79\%$ of agents that fill in payment details proceed immediately to place the order at Step 8. This structural confirmation break right before order placement blocks autonomous completion at the exact point of maximum commitment risk. To prevent operator habituation from routine prompts, this confirmation gate is triggered selectively based on transaction value, merchant category, and risk signals.

Third, high-risk commercial rails should adopt verifiable agent identity infrastructure at the protocol layer. By using decentralized identifiers and zero-knowledge verifiable credentials (such as W3C DIDs or proof-of-human bindings), a merchant or payment gateway can verify that a transacting agent is backed by an accountable human principal without exposing private identity data. While an unregistered criminal agent can evade a registry on non-compliant venues, mandatory attribution changes the incentives of ordinary deployers on regulated payment rails. This liability gives operators a clear reason to bound spending limits, restrict tool allowlists, and supervise autonomous agents carefully.
